% Source: Weinberg GR, physical PDF pp. 128--131 (printed pp. 103--106).
% Coverage: Section 4.6, Covariant Differentiation; equations
%           (4.6.1)--(4.6.19).
% Figures/tables/footnotes: none.
% Status: source-reviewed and compile-clean.
% Uncertainties: none.

\subsection{Covariant Differentiation}\label{sec:4.6}

We have already remarked that differentiation of a tensor does not generally
yield another tensor. For instance, consider a contravariant vector \(V^\mu\),
whose transformation law is
\[
  V'^\mu=\frac{\partial x'^\mu}{\partial x^\nu}V^\nu.
\]
Differentiating with respect to \(x'^\lambda\) gives
\begin{equation}
  \partial'_\lambda V'^\mu
  =
  \frac{\partial x'^\mu}{\partial x^\nu}
  \frac{\partial x^\rho}{\partial x'^\lambda}
  \partial_\rho V^\nu
  +
  \frac{\partial^2x'^\mu}
       {\partial x^\nu\,\partial x^\rho}
  \frac{\partial x^\rho}{\partial x'^\lambda}V^\nu.
  \tag{4.6.1}\label{eq:4.6.1}
\end{equation}
The first term on the right is what we would expect if
\(\partial_\lambda V^\mu\) were a tensor; the second term is what destroys
the tensor behavior.

Although \(\partial_\lambda V^\mu\) is not a tensor, we can use it to
construct a tensor. Using Eq.~\eqref{eq:4.5.8}, we see that
\begin{equation}
  \Gamma'^\mu{}_{\lambda\kappa}V'^\kappa
  =
  \frac{\partial x'^\mu}{\partial x^\nu}
  \frac{\partial x^\rho}{\partial x'^\lambda}
  \Gamma^\nu{}_{\rho\sigma}V^\sigma
  -
  \frac{\partial^2x'^\mu}
       {\partial x^\nu\,\partial x^\rho}
  \frac{\partial x^\rho}{\partial x'^\lambda}V^\nu.
  \tag{4.6.2}\label{eq:4.6.2}
\end{equation}
Adding Eqs.~\eqref{eq:4.6.1} and \eqref{eq:4.6.2}, we find that the
inhomogeneous terms cancel, yielding
\begin{equation}
  \partial'_\lambda V'^\mu
  +\Gamma'^\mu{}_{\lambda\kappa}V'^\kappa
  =
  \frac{\partial x'^\mu}{\partial x^\nu}
  \frac{\partial x^\rho}{\partial x'^\lambda}
  \left(
    \partial_\rho V^\nu+\Gamma^\nu{}_{\rho\sigma}V^\sigma
  \right).
  \tag{4.6.3}\label{eq:4.6.3}
\end{equation}
Thus we are led to define a \textit{covariant derivative}
\begin{equation}
  \nabla_\lambda V^\mu
  \equiv
  \partial_\lambda V^\mu+\Gamma^\mu{}_{\lambda\kappa}V^\kappa,
  \tag{4.6.4}\label{eq:4.6.4}
\end{equation}
and Eq.~\eqref{eq:4.6.3} tells us that \(\nabla_\lambda V^\mu\) is a
tensor:
\[
  \nabla'_\lambda V'^\mu
  =
  \frac{\partial x'^\mu}{\partial x^\nu}
  \frac{\partial x^\rho}{\partial x'^\lambda}
  \nabla_\rho V^\nu.
\]

We can also define a covariant derivative for a covariant vector \(V_\mu\).
We recall its transformation rule,
\[
  V'_\mu=\frac{\partial x^\rho}{\partial x'^\mu}V_\rho.
\]
Differentiating with respect to \(x'^\nu\) gives
\begin{equation}
  \partial'_\nu V'_\mu
  =
  \frac{\partial x^\rho}{\partial x'^\mu}
  \frac{\partial x^\sigma}{\partial x'^\nu}
  \partial_\sigma V_\rho
  +
  \frac{\partial^2x^\rho}
       {\partial x'^\mu\,\partial x'^\nu}V_\rho.
  \tag{4.6.5}\label{eq:4.6.5}
\end{equation}
From Eq.~\eqref{eq:4.5.2} we have
\begin{equation}
  \Gamma'^\lambda{}_{\mu\nu}V'_\lambda
  =
  \frac{\partial x^\rho}{\partial x'^\mu}
  \frac{\partial x^\sigma}{\partial x'^\nu}
  \Gamma^\kappa{}_{\rho\sigma}V_\kappa
  +
  \frac{\partial^2x^\kappa}
       {\partial x'^\mu\,\partial x'^\nu}V_\kappa.
  \tag{4.6.6}\label{eq:4.6.6}
\end{equation}
The inhomogeneous terms cancel if we subtract Eq.~\eqref{eq:4.6.6} from
Eq.~\eqref{eq:4.6.5}:
\begin{equation}
  \partial'_\nu V'_\mu-\Gamma'^\lambda{}_{\mu\nu}V'_\lambda
  =
  \frac{\partial x^\rho}{\partial x'^\mu}
  \frac{\partial x^\sigma}{\partial x'^\nu}
  \left(
    \partial_\sigma V_\rho-\Gamma^\kappa{}_{\rho\sigma}V_\kappa
  \right).
  \tag{4.6.7}\label{eq:4.6.7}
\end{equation}
We therefore define the covariant derivative of a covariant vector by
\begin{equation}
  \nabla_\nu V_\mu
  \equiv
  \partial_\nu V_\mu-\Gamma^\lambda{}_{\mu\nu}V_\lambda,
  \tag{4.6.8}\label{eq:4.6.8}
\end{equation}
and Eq.~\eqref{eq:4.6.7} tells us that it is a tensor:
\begin{equation}
  \nabla'_\nu V'_\mu
  =
  \frac{\partial x^\rho}{\partial x'^\mu}
  \frac{\partial x^\sigma}{\partial x'^\nu}
  \nabla_\sigma V_\rho.
  \tag{4.6.9}\label{eq:4.6.9}
\end{equation}

It is obvious how these definitions are to be extended to a general tensor.
The covariant derivative with respect to \(x^\rho\) of a tensor \(T\) equals
\(\partial_\rho T\), plus one term
\(\Gamma^\mu{}_{\rho\nu}T\) for each contravariant index \(\mu\), with
\(\mu\) replaced by \(\nu\), and minus one term
\(\Gamma^\nu{}_{\rho\lambda}T\) for each covariant index \(\lambda\), with
\(\lambda\) replaced by \(\nu\). For instance,
\begin{equation}
  \nabla_\rho T^{\mu\nu}{}_\lambda
  =
  \partial_\rho T^{\mu\nu}{}_\lambda
  +\Gamma^\mu{}_{\rho\kappa}T^{\kappa\nu}{}_\lambda
  +\Gamma^\nu{}_{\rho\kappa}T^{\mu\kappa}{}_\lambda
  -\Gamma^\kappa{}_{\rho\lambda}T^{\mu\nu}{}_\kappa.
  \tag{4.6.10}\label{eq:4.6.10}
\end{equation}
The reader may easily verify that this is a tensor.

We can also extend the idea of covariant differentiation to tensor densities.
The easiest way to do this is to recall that if
\(\mathcal T\) is a tensor density of weight \(W\), then
\((-g)^{W/2}\mathcal T\) is an ordinary tensor. Its covariant derivative is
also a tensor, and multiplying by \((-g)^{-W/2}\) gives back a tensor density
of weight \(W\). Hence the covariant derivative of a tensor density of weight
\(W\) is defined by
\begin{equation}
  \nabla_\rho\mathcal T
  \equiv
  (-g)^{-W/2}
  \nabla_\rho\!\left[(-g)^{W/2}\mathcal T\right],
  \tag{4.6.11}\label{eq:4.6.11}
\end{equation}
and we need not check that this is a tensor density of weight \(W\). The
effect is that the covariant derivative with respect to \(x^\rho\) of a
tensor density \(\mathcal T\) of weight \(W\) is constructed just as if it
were an ordinary tensor, except that we add an extra term
\((W/2g)(\partial_\rho g)\mathcal T\). For instance,
\begin{equation}
  \nabla_\rho\mathcal T^\mu{}_\lambda
  =
  \partial_\rho\mathcal T^\mu{}_\lambda
  +\Gamma^\mu{}_{\rho\nu}\mathcal T^\nu{}_\lambda
  -\Gamma^\nu{}_{\rho\lambda}\mathcal T^\mu{}_\nu
  +\frac{W}{2g}(\partial_\rho g)\mathcal T^\mu{}_\lambda.
  \tag{4.6.12}\label{eq:4.6.12}
\end{equation}

The combination of covariant differentiation with the algebraic operations
described in Section~\ref{sec:4.3} gives results similar to those for ordinary
differentiation. In particular:

\begin{enumerate}
  \item[(A)] For constant coefficients, the covariant derivative of a linear
  combination of tensors is the same linear combination of their covariant
  derivatives. For instance, if \(\alpha\) and \(\beta\) are constants, then
  \begin{equation}
    \nabla_\lambda(\alpha A^\mu+\beta B^\mu)
    =
    \alpha\nabla_\lambda A^\mu+\beta\nabla_\lambda B^\mu.
    \tag{4.6.13}\label{eq:4.6.13}
  \end{equation}

  \item[(B)] The covariant derivative of a direct product of tensors obeys
  the Leibniz rule. For instance,
  \begin{equation}
    \nabla_\rho(A^\mu B^\lambda)
    =
    (\nabla_\rho A^\mu)B^\lambda
    +A^\mu\nabla_\rho B^\lambda.
    \tag{4.6.14}\label{eq:4.6.14}
  \end{equation}

  \item[(C)] The covariant derivative of a contracted tensor is the
  contraction of the covariant derivative. For instance, setting
  \(\sigma=\lambda\) in Eq.~\eqref{eq:4.6.10} gives, after the last two
  connection terms cancel,
  \begin{equation}
    \nabla_\rho T^{\mu\lambda}{}_\lambda
    =
    \partial_\rho T^{\mu\lambda}{}_\lambda
    +\Gamma^\mu{}_{\rho\nu}T^{\nu\lambda}{}_\lambda.
    \tag{4.6.15}\label{eq:4.6.15}
  \end{equation}
\end{enumerate}

We also note that the covariant derivative of the metric tensor is zero,
because it vanishes in locally inertial coordinates where
\(\Gamma^\mu{}_{\nu\lambda}\) and \(\partial_\lambda g_{\mu\nu}\) vanish,
and a tensor that is zero in one coordinate system is zero in all systems.
The same result can be obtained more directly from Eq.~\eqref{eq:4.5.4}:
\begin{equation}
  \nabla_\lambda g_{\mu\nu}=0.
  \tag{4.6.16}\label{eq:4.6.16}
\end{equation}
(This argument can be reversed to provide yet another derivation of the
relation between \(g_{\mu\nu}\) and \(\Gamma^\lambda{}_{\mu\nu}\).) We can
also show in the same way that the covariant derivatives of the other forms
of the metric tensor vanish; that is,
\begin{equation}
  \nabla_\lambda g^{\mu\nu}=0,
  \tag{4.6.17}\label{eq:4.6.17}
\end{equation}
\begin{equation}
  \nabla_\lambda\delta^\mu{}_\nu=0.
  \tag{4.6.18}\label{eq:4.6.18}
\end{equation}
From Eqs.~\eqref{eq:4.6.16}--\eqref{eq:4.6.18} it follows that the
operations of covariant differentiation and raising and lowering indices
commute; for example,
\begin{equation}
  \nabla_\lambda(g^{\mu\nu}V_\nu)
  =
  g^{\mu\nu}\nabla_\lambda V_\nu.
  \tag{4.6.19}\label{eq:4.6.19}
\end{equation}

The importance of covariant differentiation arises from two of its
properties: It converts tensors to other tensors, and it reduces to ordinary
differentiation in the absence of gravitation, that is, when
\(\Gamma^\mu{}_{\nu\lambda}=0\). These properties suggest the following
algorithm for assessing the effects of gravitation on physical systems:

\begin{quote}
\itshape
Write the appropriate special-relativistic equations that hold in the
absence of gravitation, replace \(\eta_{\mu\nu}\) with \(g_{\mu\nu}\), and
replace all derivatives with covariant derivatives. The resulting equations
will be generally covariant and true in the absence of gravitation, and
therefore, according to the Principle of General Covariance, they will be
true in the presence of gravitational fields, provided always that we work
on a spacetime scale sufficiently small compared with the scale of the
gravitational field.
\end{quote}
